\chapter{False Vacuum Decay and Bubble Nucleation}
\label{ch:vacuum-decay}

Having established that first-order phase transitions proceed via the
nucleation of bubbles of the true vacuum within the false vacuum, we
now develop the quantitative framework for computing the nucleation
rate. The central object is the bounce solution, a saddle-point of the
Euclidean action that dominates the path integral for the false vacuum
decay rate. We begin with quantum mechanics, where the essential
physics is transparent, and then generalize to quantum field theory.

\section{Decay Rate in Quantum Mechanics}

Consider a quantum mechanical system with a metastable ground state
localized at $x_0$. If the energy eigenvalue is complex,
\begin{equation}
	E_0 = \text{Re}\,E_0 + i\,\text{Im}\,E_0,
\end{equation}
the norm of the wave function decays as
\begin{align}
	|\psi_0(t)|^2 & = \left|e^{-iHt}\psi_0(0)\right|^2 \notag      \\
	              & = e^{2\,\text{Im}\,E_0\,t}|\psi_0(0)|^2 \notag \\
	              & = e^{-\Gamma t}|\psi_0(0)|^2,
\end{align}
where $\Gamma = -2\,\text{Im}\,E_0$ is the decay rate. To compute
$E_0$, we use the path integral representation of the transition
amplitude,
\begin{equation}
	\langle x_b|e^{-i(t_b-t_a)H}|x_a\rangle
	= \int_{x(t_a)=x_a}^{x(t_b)=x_b}\mathcal{D}x\,e^{iS(x)},
	\label{eq:tr_amp}
\end{equation}
with
\begin{equation}
	S(x) = \int_{t_a}^{t_b}dt\left[
		\frac{1}{2}\left(\frac{dx}{dt}\right)^2 - V(x)\right].
\end{equation}
We proceed in three steps. First, set $|x_b\rangle = |x_a\rangle =
	|x_0\rangle$. Second, Wick rotate $t = -i\tau$ with $\tau$ real.
Third, set $\tau_a = -\hat\tau/2$ and $\tau_b = \hat\tau/2$.

The left-hand side of \eqref{eq:tr_amp} becomes
\begin{equation}
	\langle x_0|e^{-\hat H\hat\tau}|x_0\rangle
	= \sum_n e^{-E_n\hat\tau}|\langle x_0|\psi_n\rangle|^2
	\;\overset{\hat\tau\to\infty}{=}\;
	e^{-E_0\hat\tau}|\langle x_0|\psi_0\rangle|^2,
\end{equation}
where all excited states are exponentially suppressed. The ground state
energy is therefore
\begin{equation}
	E_0 = -\lim_{\hat\tau\to\infty}\frac{1}{\hat\tau}
	\log\langle x_0|e^{-\hat H\hat\tau}|x_0\rangle.
\end{equation}
The right-hand side of \eqref{eq:tr_amp} after Wick rotation becomes
\begin{equation}
	\int\mathcal{D}x\,e^{-S_E},
	\label{eq:Euclidean_integral}
\end{equation}
where the Euclidean action is
\begin{equation}
	S_E = \int_{-\hat\tau/2}^{\hat\tau/2}d\tau\left[
		\frac{1}{2}\left(\frac{dx}{d\tau}\right)^2 + V(x)\right].
	\label{eq:Euclidean_action}
\end{equation}
We evaluate \eqref{eq:Euclidean_integral} by expanding around a
stationary point $\bar x(\tau)$ of $S_E$. The Euler-Lagrange equation
gives
\begin{equation}
	\frac{d^2 x(\tau)}{d\tau^2} = V'(x),
	\label{eq:eom}
\end{equation}
with boundary condition $x(-\hat\tau/2) = x(\hat\tau/2) = x_0$.
Crucially, \eqref{eq:eom} describes a classical particle moving in the
inverted potential $-V(x)$.

Setting $x(\tau) = \bar x(\tau) + z(\tau)$ and expanding $S_E$,
\begin{align}
	S_E(x) & = S_E(\bar x)
	- \int d\tau\left(\frac{d^2\bar x}{d\tau^2}
	- V'(\bar x)\right)z \notag                    \\
	       & \quad+ \frac{1}{2}\int d\tau\,z(\tau)
	\left[-\frac{d^2}{d\tau^2} + V''(\bar x)\right]z(\tau)
	+ \mathcal{O}(z^3),
\end{align}
where integration by parts was used and boundary terms dropped. The
linear term vanishes by the equation of motion, leaving
\begin{equation}
	\int\mathcal{D}x\,e^{-S_E}
	= e^{-S_E(\bar x)}\int\mathcal{D}z\,
	\exp\left(-\frac{1}{2}\int d\tau\,z
	\left[-\frac{d^2}{d\tau^2}+V''(\bar x)\right]z\right).
\end{equation}
The Gaussian integral over $z$ gives a functional determinant,
\begin{equation}
	= N\left[\det\left(-\frac{d^2}{d\tau^2}
		+ V''(\bar x)\right)\right]^{-1/2}
	\;\overset{\hat\tau\to\infty}{=}\;
	\left(\frac{\omega}{\pi}\right)^{1/2}e^{-\omega\hat\tau/2},
\end{equation}
where $\omega = \sqrt{V''(x_0)}$. For the trivial stationary path
$\bar x(\tau) = x_0$, we have $S_E = 0$ and the transition amplitude
gives
\begin{equation}
	\langle x_0|e^{-H\hat\tau}|x_0\rangle
	= \left(\frac{\omega}{\pi}\right)^{1/2}e^{-\omega\hat\tau/2},
\end{equation}
from which $E_0 = \omega/2$.

\section{The Bounce Solution}

The non-trivial stationary path is the bounce: a solution to
\eqref{eq:eom} in the inverted potential $-V(x)$ that starts at $x_0$
at $\tau = -\infty$, rolls to the turning point $x_*$, and returns to
$x_0$ at $\tau = +\infty$.

Let $\bar x(\tau)$ be a bounce solution with $S_E(\bar x) = B$.
Its contribution to the path integral involves the correction factor
\begin{equation}
	K = \left[\frac{\det'(-\partial_\tau^2 + V''(\bar x))}
		{\det(-\partial_\tau^2 + \omega^2)}\right]^{-1/2},
\end{equation}
where the prime denotes exclusion of the zero eigenvalue.

\section{The Dilute Gas Approximation}

For $n$ widely separated bounces with centers
$\tau_1 \ll \tau_2 \ll \cdots \ll \tau_n$, integrating over all
center positions gives
\begin{equation}
	\int_{-\hat\tau/2}^{\hat\tau/2}d\tau_1
	\int_{\tau_1}^{\hat\tau/2}d\tau_2\cdots
	\int_{\tau_{n-1}}^{\hat\tau/2}d\tau_n
	= \frac{\hat\tau^n}{n!}.
\end{equation}
Summing over all $n$,
\begin{align}
	\int\mathcal{D}x\,e^{-S_E}
	 & = \left(\frac{\omega}{\pi}\right)^{1/2}e^{-\omega\hat\tau/2}
	\sum_{n=0}^{\infty}\frac{\hat\tau^n}{n!}e^{-Bn}K^n \notag       \\
	 & = \left(\frac{\omega}{\pi}\right)^{1/2}
	e^{-\hat\tau(\omega/2 - Ke^{-B})},
\end{align}
giving
\begin{equation}
	\boxed{E_0 = \frac{\omega}{2} - Ke^{-B}.}
\end{equation}
The dilute gas approximation is justified when $B \gg 1$, since the
bounce density $n/\hat\tau \sim Ke^{-B}$ is exponentially suppressed.

There are two subtleties in computing $K$. First, the operator
$(-\partial_\tau^2 + V''(\bar x))$ has a zero eigenvalue from
translational invariance of the bounce, which must be excluded from
the determinant and treated separately, contributing $(B/2\pi)^{1/2}$.
Second, it has a negative eigenvalue requiring analytic continuation.
Following Callan and Coleman, the decay rate per unit time is
\begin{equation}
	\Gamma = \left(\frac{B}{2\pi}\right)^{1/2}
	\left[\frac{\det'(-\partial_\tau^2+V''(\bar x))}
		{\det(-\partial_\tau^2+\omega^2)}\right]^{-1/2}e^{-B}.
\end{equation}

\section{Decay Rate in Quantum Field Theory}

The generalization to QFT proceeds identically. The decay rate per
unit time and volume is
\begin{equation}
	\Gamma = \frac{B^2}{4\pi^2}
	\left[\frac{\det'(-\partial_\tau^2-\nabla^2+V''(\bar\phi))}
		{\det(-\partial_\tau^2-\nabla^2+\omega^2)}\right]^{-1/2}e^{-B},
\end{equation}
where the bounce action is
\begin{equation}
	B = S_E[\bar\phi] = \int d\tau\,d^3x\left[
		\frac{1}{2}(\partial_\tau\bar\phi)^2
		+ \frac{1}{2}(\nabla\bar\phi)^2
		+ V(\bar\phi)\right].
\end{equation}
The functional determinant contributes only a polynomial prefactor;
the exponential $e^{-B}$ dominates.

The bounce $\bar\phi$ satisfies
\begin{equation}
	(\partial_\tau^2 + \nabla^2)\bar\phi = V'(\bar\phi),
	\label{eq:bounce_eom}
\end{equation}
with boundary conditions $\bar\phi \to \phi_0$ as
$\tau \to \pm\infty$ and $|\mathbf{x}| \to \infty$. By a theorem of
Coleman, Glashow, and Martin (1978), the $O(4)$-symmetric solution
has the lowest action. Writing $\bar\phi = \phi(\rho)$ with
$\rho = \sqrt{\tau^2 + |\mathbf{x}|^2}$, the equation reduces to
\begin{equation}
	\frac{d^2\bar\phi}{d\rho^2}
	+ \frac{3}{\rho}\frac{d\bar\phi}{d\rho} = V'(\bar\phi),
	\label{eq:spherical_eom}
\end{equation}
with $\bar\phi'(0) = 0$ and $\bar\phi \to \phi_0$ as
$\rho \to \infty$.

\section{The Overshoot-Undershoot Algorithm}

Equation \eqref{eq:spherical_eom} is solved numerically by starting
from $\bar\phi(0) = \phi_i \in [\phi_0, \phi_1]$ with
$\bar\phi'(0) = 0$ and evolving to large $\rho$. If $\phi_i$ is too
large the field overshoots past $\phi_0$; if too small it undershoots
and never reaches $\phi_0$. The correct bounce is found by binary
search. The two limiting cases are the thin-wall approximation, where
the potential barrier is narrow compared to the bubble radius, and the
thick-wall approximation, where it is wide.

\section{Bubble Nucleation and Dynamics}

The bounce has a dual role: it dominates the false vacuum decay rate,
and it describes the field configuration after tunneling. The initial
bubble is
\begin{align}
	\phi_{\text{bubble}}(t=0,\mathbf{x})
	                                          & = \bar\phi(\tau=0,\mathbf{x}), \\
	\partial_t\phi_{\text{bubble}}\big|_{t=0} & = 0.
\end{align}
After nucleation the bubble evolves by analytic continuation from
Euclidean space,
\begin{equation}
	\phi_{\text{bubble}}(t,\mathbf{x})
	= \bar\phi\left(\sqrt{|\mathbf{x}|^2 - t^2}\right),
\end{equation}
and the bubble wall position satisfies
$|\mathbf{x}| = \sqrt{t^2 + C^2}$, approaching the speed of light at
late times.

\section{Finite Temperature and the O(3) Bounce}

At finite temperature the Euclidean time direction is compactified to
a circle of size $\beta = T^{-1}$. At high temperature the
$O(4)$-symmetric bounce transitions to an $O(3)$-symmetric,
time-independent bounce (Linde, 1983). Setting
$\partial_\tau\bar\phi = 0$, the Euclidean action reduces to
\begin{equation}
	S_E = \frac{1}{T}\int d^3x\left[
		\frac{1}{2}(\nabla\bar\phi)^2 + V(\bar\phi)\right]
	\equiv \frac{S_3}{T}.
\end{equation}
The two cases are unified as
\begin{equation}
	\frac{d^2\phi}{d\rho^2}
	+ \frac{c_i}{\rho}\frac{d\phi}{d\rho} = V'(\phi),
	\qquad c_i = \begin{cases}
		3 & O(4)\text{ symmetry}, \\
		2 & O(3)\text{ symmetry},
	\end{cases}
\end{equation}
with bounce actions
\begin{equation}
	B = \begin{cases}
		S_E[\bar\phi] & O(4), \\
		S_3/T         & O(3).
	\end{cases}
\end{equation}
At finite temperature, $V(\bar\phi)$ is replaced by
$V_{\text{eff}}^\beta(\bar\phi)$ and the nucleation rate is
$\Gamma \sim e^{-S_3/T}$. The nucleation temperature $T_n$ and
percolation temperature $T_p$, together with the latent heat $\alpha$
and transition rate $\tilde\beta$, characterize the transition and
determine the gravitational wave spectrum produced during bubble
collisions.
